% CP-VI-0025
% target_chapter: VI/15 — Kernels, Images, and Quotients
% source_unit_id: IMP-DL-A99C7658F3-P001
% source: Downloads/theory-of-algebraic-geometry-2.tex L183-L351
% solution_provenance: SOURCE_EXPLICIT_SOLUTION

\subsection*{Problem 7. A Morphism of Sheaves is an Isomorphism if and only if it is Injective and Surjective}

Let
\[
\varphi:\mathcal{F} \longrightarrow \mathcal{G}
\]
be a morphism of sheaves.

\textbf{Claim.} The morphism \(\varphi\) is an isomorphism if and only if it is injective and surjective.

\textbf{Solution.}

First suppose that \(\varphi\) is an isomorphism of sheaves. Then there exists a morphism
\[
\psi:\mathcal{G}\longrightarrow \mathcal{F}
\]
such that
\[
\psi\circ \varphi=\operatorname{id}_{\mathcal{F}},
\qquad
\varphi\circ \psi=\operatorname{id}_{\mathcal{G}}.
\]
Passing to stalks, for every \(x\in X\), we obtain maps
\[
\varphi_x:\mathcal{F}_x\longrightarrow \mathcal{G}_x,
\qquad
\psi_x:\mathcal{G}_x\longrightarrow \mathcal{F}_x
\]
such that
\[
\psi_x\circ \varphi_x=\operatorname{id}_{\mathcal{F}_x},
\qquad
\varphi_x\circ \psi_x=\operatorname{id}_{\mathcal{G}_x}.
\]
Hence \(\varphi_x\) is an isomorphism of abelian groups for every \(x\in X\). Therefore \(\varphi_x\) is injective and surjective for every \(x\). Thus \(\varphi\) is injective and surjective.

Conversely, suppose that \(\varphi\) is injective and surjective. This means that for every \(x\in X\), the induced morphism on stalks
\[
\varphi_x:\mathcal{F}_x\longrightarrow \mathcal{G}_x
\]
is injective and surjective. Hence \(\varphi_x\) is bijective for every \(x\in X\).

We will show that \(\varphi\) is an isomorphism of sheaves.

Let \(U\subseteq X\) be open and let
\[
t\in \mathcal{G}(U).
\]
We want to construct a unique section
\[
s\in \mathcal{F}(U)
\]
such that
\[
\varphi_U(s)=t.
\]

Since \(\varphi_x:\mathcal{F}_x\to \mathcal{G}_x\) is surjective for every \(x\in U\), for each point \(x\in U\), the germ
\[
t_x\in \mathcal{G}_x
\]
has a preimage in \(\mathcal{F}_x\). Therefore there exists a germ
\[
s_x\in \mathcal{F}_x
\]
such that
\[
\varphi_x(s_x)=t_x.
\]

By the definition of a germ, for each \(x\in U\), there exists an open neighbourhood
\[
V_x\subseteq U
\]
and a section
\[
s_{V_x}\in \mathcal{F}(V_x)
\]
representing the germ \(s_x\). Since
\[
\varphi_x(s_x)=t_x,
\]
after shrinking \(V_x\) if necessary, we may assume that
\[
\varphi(s_{V_x})=t|_{V_x}.
\]

Now consider two such neighbourhoods \(V_x\) and \(V_y\). On the overlap \(V_x\cap V_y\), we have
\[
\varphi(s_{V_x}|_{V_x\cap V_y})
=
t|_{V_x\cap V_y}
=
\varphi(s_{V_y}|_{V_x\cap V_y}).
\]
Since \(\varphi\) is injective, the map
\[
\mathcal{F}(V_x\cap V_y)\longrightarrow \mathcal{G}(V_x\cap V_y)
\]
is injective. Hence
\[
s_{V_x}|_{V_x\cap V_y}
=
s_{V_y}|_{V_x\cap V_y}.
\]

Thus the local sections \(s_{V_x}\) agree on overlaps. By the gluing axiom for the sheaf \(\mathcal{F}\), they glue to a unique global section
\[
s\in \mathcal{F}(U)
\]
such that
\[
s|_{V_x}=s_{V_x}
\]
for every \(x\in U\).

Then
\[
\varphi(s)|_{V_x}
=
\varphi(s|_{V_x})
=
\varphi(s_{V_x})
=
t|_{V_x}.
\]
Since these equalities hold on an open cover of \(U\), and since \(\mathcal{G}\) is a sheaf, it follows that
\[
\varphi(s)=t.
\]
Therefore
\[
\varphi_U:\mathcal{F}(U)\longrightarrow \mathcal{G}(U)
\]
is surjective for every open set \(U\).

Now let \(s\in \mathcal{F}(U)\) and suppose
\[
\varphi(s)=0.
\]
Then for every \(x\in U\),
\[
\varphi_x(s_x)=0.
\]
Since \(\varphi_x\) is injective, it follows that
\[
s_x=0
\]
for every \(x\in U\). A section of a sheaf is zero if and only if all of its germs are zero. Hence
\[
s=0.
\]
Therefore
\[
\varphi_U:\mathcal{F}(U)\longrightarrow \mathcal{G}(U)
\]
is injective for every open set \(U\).

Thus \(\varphi_U\) is bijective for every open set \(U\). The inverse maps are compatible with restrictions because \(\varphi\) is a morphism of sheaves. Hence \(\varphi\) is an isomorphism of sheaves.

\[
\boxed{
	\varphi \text{ is an isomorphism }
	\Longleftrightarrow
	\varphi \text{ is injective and surjective.}
}
\]
